\section{数据基础与潜在数据需求}

\subsection{现有数据基础}

结合当前已整理的数据资料，本研究已经具备开展西安轨道交通系统韧性分析的初步数据基础，主要包括以下几类：

\begin{enumerate}
\item \textbf{轨道交通刷卡交易数据} 地铁票+长安通刷卡数据，可用提取乘客进出站时间、进出站站点以及时段客流分布特征
\item \textbf{轨道交通网络与站点空间数据} 当前已整理西安轨道交通站点主表、线路信息、区间网络结构以及站点空间坐标数据，可用于构建站点--区间--换乘节点构成的轨道交通网络，为后续开展网络韧性测度与空间分布分析提供基础
\item \textbf{站点周边 POI 与站点环境数据} 基于高德数据整理站点周边一定缓冲区内的 POI 原始数据、汇总数据及站点缓冲区空间范围，可用于表征站点周边功能结构、活动吸引类型与建成环境差异
\end{enumerate}


\subsection{潜在数据需求}

\begin{enumerate}
\item \textbf{天气观测与天气扰动数据} 
\item \textbf{站点设施与运行组织数据} 站点层面的设施属性与运行组织信息，例如换乘方式、出入口数量、站台规模、列车发车间隔、断面运能及限流措施等
\item \textbf{更丰富的站点环境与建成环境指标}可进一步补充土地利用类型、人口密度、就业密度、道路网络可达性、公交接驳条件以及重要公共服务设施分布等变量
\item \textbf{日历与事件数据。} 节假日、工作日/周末校放假时间以及大型活动日程等辅助数据
\end{enumerate}


